\documentclass[11pt]{article}
\usepackage[T1]{fontenc}
\usepackage{lmodern}
\usepackage[margin=1in,headheight=15pt]{geometry}
\usepackage{amsmath,amssymb,amsthm,mathtools}
\usepackage{booktabs,array,longtable,microtype,xurl,fancyhdr}
\usepackage[dvipsnames]{xcolor}
\usepackage[colorlinks=true,linkcolor=MidnightBlue,citecolor=MidnightBlue,urlcolor=MidnightBlue]{hyperref}
\usepackage{bookmark}
\emergencystretch=3em
\setlength{\parskip}{3pt}
\numberwithin{equation}{section}
\newtheorem{theorem}{Theorem}[section]
\newtheorem{lemma}[theorem]{Lemma}
\newtheorem{proposition}[theorem]{Proposition}
\newtheorem{corollary}[theorem]{Corollary}
\newtheorem{conjecture}[theorem]{Conjecture}
\theoremstyle{definition}\newtheorem{definition}[theorem]{Definition}
\theoremstyle{remark}\newtheorem{remark}[theorem]{Remark}
\newcommand{\dd}{\,\mathrm d}
\newcommand{\ii}{\mathrm i}
\newcommand{\R}{\mathbb R}
\newcommand{\C}{\mathbb C}
\newcommand{\E}{\mathbb E}
\newcommand{\cut}{\C\setminus[1,\infty)}
\newcommand{\Cstar}{C_*}
\DeclareMathOperator{\Li}{Li}
\DeclareMathOperator{\Cov}{Cov}
\DeclareMathOperator{\Impart}{Im}
\DeclareMathOperator{\Repart}{Re}
\DeclareMathOperator{\Beta}{Beta}
\DeclareMathOperator{\GammaLaw}{Gamma}
\newcommand{\basecommit}{\nolinkurl{a0a90ef31877f98be437191c48b46f02d5456867}}
\pagestyle{fancy}\fancyhf{}
\fancyhead[L]{\small Critical Euler errors for harmonic polylogarithms}
\fancyhead[R]{\small ProveIt}
\fancyfoot[C]{\thepage}
\hypersetup{pdftitle={Critical Euler Errors for Harmonic Polylogarithms},pdfauthor={Research continuation prepared for Vladimir Reshetnikov},pdfsubject={Sharp constants, positive surplus kernels, and a drifting extremizer}}
\title{\textbf{Critical Euler Errors\\for Harmonic Polylogarithms}\\[7pt]
\large Sharp constants, positive surplus kernels, and a drifting extremizer}
\author{Research continuation prepared for Vladimir Reshetnikov\\
\small Prepared with OpenAI assistance for the ProveIt programme}
\date{9 October 2026 (America/Los\_Angeles)}
\begin{document}
\maketitle\thispagestyle{empty}
\begin{abstract}
Let $F_{a,b}(z)=\sum_{n\ge1}H_{n-1}^{(b)}z^n/n^a$, with principal
continuation to $\cut$, and let $g_{a,b}=\Impart F_{a,b}(\ii)$.
We prove the sharp-critical-constant conjecture in the incoming report
\emph{The Sharp Real-Order Threshold for Harmonic Polylogarithms}:
$a\mapsto-2g_{a,1-a}$ is strictly decreasing on $(0,1)$, and the optimal
uniform constant in its one-sided Euler bound is
$\Cstar=\pi/4+(\log2)/2$.
The proof comes from a fixed-total-order beta--gamma transport. It gives
strict order monotonicity for every total order $a+b>0$ at imaginary
radii up to $\sqrt3$, as well as exact endpoint comparisons on the real
slit interval.

A second argument proves that the critical negative density is strictly
larger than Lebesgue density. The resulting positive surplus has an exact
balanced Bose-kernel representation, strengthens the harmonic-defect
Hausdorff moment inequalities, and yields a zero-free Stieltjes
factorization for a difference of polylogarithms.

For the scaled Euler error $R_N(a)=2^N(E_N-g_{a,1-a})$, we establish a
uniform asymptotic and determine its moving maximizer. With
$L=\frac12\log N$ and $\ell=\log L$, every maximizing outer order $a_N$
satisfies $(1-a_N)\ell\to1$, while
$\max_aR_N(a)=\frac{\sqrt\pi}{2\sqrt N}
[1+(1+o(1))/(eL\ell)]$.
A refinement identifies a coefficient involving the first Stieltjes
constant. Exact rational certificates demonstrate the finite-depth
reversal $R_{128}(1/2)>R_{128}(1/4)$.
Thus the proved monotonicity of the sharp all-depth constant does not
extend to each individual scaled error. All analytic assertions are
proved in the text; exact arithmetic, numerical diagnostics, and remaining
conjectures are explicitly separated.
\end{abstract}
\clearpage
\begingroup\small\setlength{\parskip}{0pt}\tableofcontents\endgroup
\clearpage

\section{Scope, baseline, and the question resolved}
\label{sec:scope}
The canonical ProveIt manuscript \emph{Polylogarithms and their Arithmetic
Bridges} develops signed kernels and certified evaluation alongside
cyclotomic reductions and arithmetic special values \cite{manuscript}.
The incoming report \emph{The Sharp Real-Order Threshold for Harmonic
Polylogarithms} extends the finite signed moment representation to the
sharp domain $a+b\ge1$ and identifies the required atom on its boundary
\cite{threshold}. Its Conjecture~12.3, labelled
\texttt{conj:constant} in the delivered source, asks for the strict decrease
of $-2g_{a,1-a}$ and the optimal uniform Euler constant. We settle that
conjecture and then examine an optimization problem it does not settle:
which parameter is worst at a prescribed finite Euler depth?

The repository revision inspected for this delivery is
\begin{center}\small\basecommit.\end{center}
The incoming archive \texttt{ProveIt\_RealOrder\_Threshold.zip} was read
in its delivered form; the accessible copy agrees with the incoming
repository blob. The precise source path, archive identity, and audit
boundary are recorded in \texttt{data/provenance.json}. This is a targeted
continuation, not a line-by-line audit of every incoming archive or every
chapter. No upstream file or branch has been changed.

The finite-measure threshold, the basic double kernel, critical angular
uniqueness, and the fixed-parameter Euler remainder formula below are
\emph{prior results}. We include short derivations of the identities needed
here so that the new arguments can be checked independently. The new
results are the fixed-total-order comparisons, proof of the named
conjecture, strict surplus density, strengthened moment consequences,
uniform critical error expansion, and extremizer drift. We do not claim
global literature priority for every identity, and none of the analytic
proofs has been formalized in a proof assistant or independently refereed.
The classical beta, gamma, and zeta facts are cited at their point of use.

\subsection{Notation and summation conventions}
For real $a,b>0$, define
\begin{equation}\label{eq:def}
 c_n(a,b)=\frac{H_{n-1}^{(b)}}{n^a},\qquad
 H_m^{(b)}=\sum_{j=1}^m j^{-b},\qquad
 F_{a,b}(z)=\sum_{n=1}^\infty c_n(a,b)z^n\quad(|z|<1).
\end{equation}
Thus $c_1=0$ and $c_2=2^{-a}$. At positive integer orders,
$F_{a,b}(z)=\Li_{a,b}(z,1)$ in the decreasing-index convention.
We write
\[
 G_{a,b}(\rho)=-\Impart F_{a,b}(\ii\rho),\qquad
 g_{a,b}=\Impart F_{a,b}(\ii),\qquad w=a+b.
\]
The principal polylogarithm conventions are those of DLMF~25.12
\cite{dlmf-polylog}. On $a+b=1$ the coefficients tend to $1/a$, so the
ordinary boundary series fails the term test. All Gaussian values on
that line mean the analytic continuation, equivalently the Abel value.
Euler transformation below is justified by a measure identity, not by
pretending the original boundary series converges.

Unless stated otherwise, a constant implicit in $O(\cdot)$ may depend on
fixed parameters. The word \emph{uniform} always identifies the parameter
set over which the constant and the threshold are independent. We use
$\Delta Q_n=Q_n-Q_{n+1}$ for the decreasing finite-difference convention.

\section{Fixed-total-order beta--gamma transport}
\label{sec:transport}
\subsection{A positive double kernel and a different coupling}
The useful part of the earlier double integral is not merely its positivity.
A change of variables makes the total order, rather than the two separate
orders, the gamma shape.

\begin{lemma}[Double kernel, recalled with proof]\label{lem:double}
For $a,b>0$ and $z\in\cut$,
\begin{equation}\label{eq:double}
 F_{a,b}(z)=z^2\int_{0<y<x<1}
 \frac{(-\log x)^{a-1}(\log(x/y))^{b-1}}
 {\Gamma(a)\Gamma(b)(1-zx)(1-zy)}\dd y\dd x.
\end{equation}
The positive weight, without the two denominator factors, has mass $2^{-a}$.
\end{lemma}
\begin{proof}
Expand the denominators for $|z|<1$. For nonnegative integers $j,k$,
substitution $y=xv$ and two gamma integrals give
\[
 \int_{0<y<x<1}
 \frac{(-\log x)^{a-1}(\log(x/y))^{b-1}}
 {\Gamma(a)\Gamma(b)}x^jy^k\dd y\dd x
 =\frac1{(j+k+2)^a(k+1)^b}.
\]
Summation over $j+k=n-2$ yields $c_n$. Absolute convergence justifies the
expansion, and $j=k=0$ gives the mass. On compact subsets of the slit plane
both denominators stay uniformly away from zero. Dominated holomorphic
integration then continues the identity to the claimed domain.
\end{proof}

\begin{theorem}[Fixed-weight transport identity]\label{thm:transport}
Let $T\sim\GammaLaw(w,1)$ and $V\sim\Beta(a,b)$ be independent, where
$w=a+b$. Put $X=e^{-TV}$ and $Y=e^{-T}$. For $z\in\cut$,
\begin{equation}\label{eq:transport}
 \frac{F_{a,b}(z)}{z^2}
 =\E\left[\frac{X}{(1-zX)(1-zY)}\right],
\end{equation}
with the quotient continuously extended at zero. In particular,
\begin{equation}\label{eq:Gtransport}
 G_{a,b}(\rho)=\rho^3\E[H_\rho(X,Y)],\qquad
 H_\rho(x,y)=\frac{x(x+y)}{(1+\rho^2x^2)(1+\rho^2y^2)}.
\end{equation}
\end{theorem}
\begin{proof}
In \eqref{eq:double}, set $U=-\log x$, $S=\log(x/y)$,
$T=U+S$, and $V=U/T$. Then $x=e^{-TV}$, $y=e^{-T}$, and
\[
 \dd x\dd y=e^{-2U-S}\dd U\dd S,
 \qquad \dd U\dd S=T\dd T\dd V.
\]
Since $e^{-2U-S}=e^{-T}e^{-TV}$, the resulting measure is $X$ times
\[
 \frac{e^{-T}T^{w-1}}{\Gamma(w)}\dd T\;
 \frac{V^{a-1}(1-V)^{b-1}}{\mathrm B(a,b)}\dd V.
\]
Euler's beta identity $\mathrm B(a,b)=\Gamma(a)\Gamma(b)/\Gamma(a+b)$
\cite{dlmf-beta} proves \eqref{eq:transport}. The imaginary part at
$z=\ii\rho$ is obtained by rationalizing the two denominator factors.
\end{proof}

\begin{remark}[The extra factor is essential]
The expectation in \eqref{eq:transport} has the numerator $X$. Omitting
it would produce the wrong coefficients and invalidate the monotonicity
argument. The earlier normalization using independent gamma variables at
rates two and one instead has a prefactor $2^{-a}$. The coupling here
uses rate one for the common total $T$ and keeps the extra factor inside
the expectation. This is a change of representation, not a correction
to the earlier correctly normalized formula.
\end{remark}

\subsection{Strict stochastic comparison}
\begin{lemma}[Beta score covariance]\label{lem:beta}
Fix $w>0$. If $V_a\sim\Beta(a,w-a)$, $0<a<w$, and $f$ is bounded and
strictly decreasing on $(0,1)$, then $\E f(V_a)$ is strictly decreasing
in $a$. On compact subintervals of $(0,w)$,
\begin{equation}\label{eq:covariance}
 \frac{\dd}{\dd a}\E f(V_a)
 =\Cov\left(f(V_a),\log\frac{V_a}{1-V_a}\right)<0.
\end{equation}
The derivative formula also holds for bounded measurable $f$ without a
strictness assertion.
\end{lemma}
\begin{proof}
Differentiate the normalized beta density. Its score is the centered
version of $\log(v/(1-v))$, so differentiation gives
\eqref{eq:covariance}. The logarithmic factors are integrable uniformly
on compact shape-parameter intervals. If $V_a'$ is an independent copy,
then the covariance equals
\[
 \frac12\E\left[(f(V_a)-f(V_a'))
 \left(\log\frac{V_a}{1-V_a}-\log\frac{V_a'}{1-V_a'}\right)\right].
\]
The product is strictly negative whenever the two variables differ.
Their densities are positive and continuous in the interior, so equality
has probability zero. This proves strictness.
\end{proof}

\begin{theorem}[Fixed-total-order comparisons]\label{thm:monotonicity}
For every $w>0$ and $0<\rho\le\sqrt3$, the function
\begin{equation}\label{eq:orderdecrease}
 a\longmapsto G_{a,w-a}(\rho),\qquad 0<a<w,
\end{equation}
is strictly decreasing. For every real $z<1$, the function
$a\mapsto F_{a,w-a}(z)/z^2$ is also strictly decreasing, with its
continuous interpretation at $z=0$.
\end{theorem}
\begin{proof}
For $0<y<x<1$, direct differentiation gives
\begin{equation}\label{eq:Hderivative}
 \partial_xH_\rho(x,y)=
 \frac{2x+y-\rho^2x^2y}
 {(1+\rho^2y^2)(1+\rho^2x^2)^2}>0.
\end{equation}
Indeed, for $\rho^2\le3$ the numerator is at least
$2x+y-3x^2y\ge3y(1-x^2)>0$.
For fixed $T>0$, $Y=e^{-T}$ is fixed and $X=e^{-TV}$ decreases strictly
with $V$. Thus $H_\rho(e^{-TV},e^{-T})$ is bounded and strictly
decreasing in $V$. Apply Lemma~\ref{lem:beta} conditionally on $T$ and
then integrate its strict inequality against the common gamma law.
This proves \eqref{eq:orderdecrease}.

For real $z<1$, replace $H_\rho$ by
$x/((1-zx)(1-zy))$. Its $x$-derivative is
\[
 \frac1{(1-zy)(1-zx)^2}>0.
\]
The same proof applies. At $z=0$, the quotient is $2^{-a}$, consistently
with the continuous extension. All integrands used here are uniformly
bounded for the fixed values of $z$ or $\rho$ in the theorem.
\end{proof}

The radius $\sqrt3$ is a sufficient domain obtained from the pointwise
kernel inequality. We do not assert that it is the maximal domain of
integrated order monotonicity. When $\rho>\sqrt3$, the numerator in
\eqref{eq:Hderivative} can be negative near $x=y=1$, but that fact alone
is not a counterexample for the expectation.

\subsection{Endpoint identities and explicit Gaussian bounds}
\begin{proposition}[Order-zero endpoints]\label{prop:endpoints}
For fixed $w>0$, locally uniformly for $z\in\cut$,
\begin{align}
 \lim_{a\downarrow0}F_{a,w-a}(z)
 &=\frac{z}{1-z}\Li_w(z),\label{eq:endleft}\\
 \lim_{a\uparrow w}F_{a,w-a}(z)
 &=\Li_{w-1}(z)-\Li_w(z).\label{eq:endright}
\end{align}
\end{proposition}
\begin{proof}
The beta variable tends in probability to zero as $a\downarrow0$ and
to one as $a\uparrow w$: its respective means or complementary means
tend to zero. Bounded convergence in \eqref{eq:transport} therefore
supplies both limiting holomorphic functions, uniformly on compact
slit sets. In the disk their coefficients are respectively
$H_{n-1}^{(w)}$ and $(n-1)n^{-w}$. Summing the first sequence by a
geometric series and splitting the second proves the displayed formulas.
\end{proof}

Define $\beta(s)=\Impart\Li_s(\ii)$ and
$\eta(s)=-\Li_s(-1)$ by continuation for real $s$.
For positive $w$ the second quantity is the usual Dirichlet eta value.
The parity split in the absolutely convergent disk series, continued to
$z=\ii$, gives
$\Repart\Li_w(\ii)=-2^{-w}\eta(w)$.
Consequently Theorem~\ref{thm:monotonicity} implies the strict bounds
\begin{equation}\label{eq:Gaussianbounds}
 \beta(w)-\beta(w-1)
 <-g_{a,b}<\frac{\beta(w)+2^{-w}\eta(w)}2,
 \qquad a,b>0,\quad a+b=w.
\end{equation}
Both endpoints are sharp as limiting parameter values. On the real
interval $z<1$, the analogous strict comparison is
\begin{equation}\label{eq:realbounds}
 \frac{\Li_{w-1}(z)-\Li_w(z)}{z^2}
 <\frac{F_{a,b}(z)}{z^2}
 <\frac{\Li_w(z)}{z(1-z)},
\end{equation}
with continuous values at zero. These are exact analytic inequalities,
not numerical integer-relation discoveries.

\section{The critical measure and the sharp Euler constant}
\label{sec:constant}
\subsection{The critical representation}
From now on $0<a<1$, $b=1-a$. Define
\begin{equation}\label{eq:hka}
 k_a(T)=\frac{T^{a-1}}{\Gamma(a)},\qquad
 h(t)=\frac1{1-e^{-t}}-\frac1t.
\end{equation}
The function $h$ increases from $1/2$ to $1$, since
$h'(t)=t^{-2}-[4\sinh^2(t/2)]^{-1}>0$.
Put
\begin{align}
 C_a(T)&=\frac1{\Gamma(a)\Gamma(b)}
       \int_0^T(T-t)^{a-1}t^{b-1}h(t)\dd t,
       \label{eq:C}\\
 p_a(T)&=-\zeta(b)k_a(T)+C_a(T).
       \label{eq:p}
\end{align}
Here $C_a(T)=\E h(TV)$ for $V\sim\Beta(b,a)$; note that the order of
these beta parameters differs from the transport variable in
Section~\ref{sec:transport}. In particular $1/2<C_a(T)<1$.

\begin{lemma}[Critical signed measure, prior identity]\label{lem:measure}
The density $p_a(-\log u)$ is positive, has mass $1/a$ on $(0,1)$,
and
\begin{equation}\label{eq:mu}
 \mu_a=\frac1a\delta_1-p_a(-\log u)\dd u
\end{equation}
has moments $\int u^{n-1}\dd\mu_a=c_n(a,1-a)$. Thus
\begin{equation}\label{eq:Stieltjes}
 F_{a,1-a}(z)=\int_{[0,1]}\frac{z}{1-zu}\dd\mu_a(u)
 \quad(z\in\cut).
\end{equation}
\end{lemma}
\begin{proof}
The regularized Hurwitz integral \cite{dlmf-hurwitz} is
\[
 \zeta(b,q)=\frac{q^{1-b}}{b-1}
 +\frac1{\Gamma(b)}\int_0^\infty e^{-qt}t^{b-1}h(t)\dd t
 \qquad(0<b<1,\ q>0).
\]
It follows by splitting the standard integral for $b>1$ and continuing
the bounded-regularizer integral to $\Re b>0$.
Taking the Laplace transform of \eqref{eq:C}, and using $a+b=1$,
gives
\begin{equation}\label{eq:plaplace}
 \int_0^\infty e^{-qT}p_a(T)\dd T
 =\frac1a-q^{-a}\bigl(\zeta(b)-\zeta(b,q)\bigr).
\end{equation}
The finite Hurwitz difference at $q=n$ is $H_{n-1}^{(b)}$.
The case $n=1$ gives the mass $1/a$. Positivity follows from
$\zeta(b)<0$, $k_a>0$, and $C_a>0$. One may obtain the zeta sign from
its positive alternating eta series and $1-2^{1-b}<0$.
Expansion of the resolvent in the disk and dominated continuation
prove \eqref{eq:Stieltjes}. This reproduces the critical part of
\cite[Theorem~3.1]{threshold}, not a new threshold assertion.
\end{proof}

\subsection{Euler errors and the conjecture}
Define, exactly as in the incoming report,
\begin{equation}\label{eq:eulerdef}
 f(n)=c_{2n+1}(a,1-a),\qquad
 d_k=\sum_{n=0}^k(-1)^n\binom kn f(n),\qquad
 E_N=\sum_{k=0}^{N-1}\frac{d_k}{2^{k+1}}.
\end{equation}
For $N\ge1$, let
\begin{equation}\label{eq:Rdef}
 R_N(a)=2^N(E_N-g_{a,1-a}),\qquad
 q_N(u)=\frac{(1-u^2)^N}{1+u^2}.
\end{equation}
The binomial theorem, followed by a geometric summation under the finite
signed measure, gives
\begin{equation}\label{eq:Rintegral}
 R_N(a)=\int_0^1q_N(u)p_a(-\log u)\dd u>0.
\end{equation}
The endpoint atom contributes zero because $N\ge1$. For fixed $a$,
these integrals decrease strictly to zero, and
\begin{equation}\label{eq:R1}
 R_1(a)=-2g_{a,1-a},\qquad
 0<E_N-g_{a,1-a}<\frac1a\,2^{-N}.
\end{equation}
These facts were already established in \cite[Section~9]{threshold}.

\begin{theorem}[Resolution of the sharp-critical-constant conjecture]
\label{thm:sharp}
The map $a\mapsto R_1(a)=-2g_{a,1-a}$ decreases strictly on $(0,1)$.
Its endpoint limits are
\begin{equation}\label{eq:criticalends}
 R_1(0)=\frac\pi4+\frac{\log2}{2}=\Cstar,
 \qquad R_1(1)=\frac\pi2-1.
\end{equation}
For every $0<a<1$ and every integer $N\ge1$,
\begin{equation}\label{eq:sharpbound}
 0<E_N-g_{a,1-a}<\Cstar\,2^{-N}.
\end{equation}
No smaller constant works uniformly over this open parameter interval
and all such $N$.
\end{theorem}
\begin{proof}
Set $w=\rho=1$ in Theorem~\ref{thm:monotonicity}.
At the two order endpoints use \eqref{eq:endleft}--\eqref{eq:endright},
$\Li_1(z)=-\log(1-z)$, and $\Li_0(z)=z/(1-z)$.
Equivalently, \eqref{eq:Gaussianbounds} uses
$\beta(1)=\pi/4$, $\beta(0)=1/2$, and $\eta(1)=\log2$.
This proves \eqref{eq:criticalends} and strict decrease.
Since $R_N(a)\le R_1(a)<\Cstar$, \eqref{eq:sharpbound} follows.
If $C<\Cstar$, choose $a$ sufficiently close to zero that $R_1(a)>C$;
then $N=1$ contradicts the proposed bound. Thus the constant is sharp.
\end{proof}

Numerically $\Cstar=1.13197175367742096432427690655\ldots$.
The theorem proves both components of the prior conjecture, not just an
upper estimate. The distinction between a strict uniform bound on an
open parameter domain and an attained maximum matters: $\Cstar$ is a
supremum at $a=0,N=1$, not a maximum with $0<a<1$.
For one fixed interior parameter, the smallest constant in a
\emph{non-strict} all-$N$ bound is $R_1(a)$ and is attained at $N=1$.

\section{A strictly positive surplus over Lebesgue density}
\label{sec:surplus}
The critical density contains a large cancellation when written as
$-\zeta(b)k_a+C_a$. Its residual above one has an unexpectedly clean
sign. This sign will be the source of stronger moment inequalities and
a sharp lower envelope for Euler errors.

\begin{lemma}[Critical power balance]\label{lem:balance}
For $0<a<1$, $b=1-a$, and $T>0$,
\begin{equation}\label{eq:balance}
 \int_0^T[k_a(T-t)-k_a(T)]t^{b-2}\dd t
 =k_a(T)\int_T^\infty t^{b-2}\dd t
 =\frac{k_a(T)T^{-a}}a.
\end{equation}
\end{lemma}
\begin{proof}
After $t=Tv$, the required scalar integral is
\[
 \int_0^1\bigl[(1-v)^{a-1}-1\bigr]v^{-a-1}\dd v=\frac1a.
\]
An antiderivative is
$-a^{-1}[(1-v)^a-1]v^{-a}$. Its lower limit is zero because
$[(1-v)^a-1]v^{-a}=O(v^{1-a})$; its upper limit is $1/a$.
Both endpoint singularities of the integrand are integrable.
Rescaling proves the identity.
\end{proof}

Set $\phi(t)=t/(e^t-1)$, which decreases strictly on $(0,\infty)$.
Its derivative has numerator $e^t(1-t)-1<0$: that numerator vanishes
at zero and has derivative $-te^t<0$.

\begin{theorem}[Positive surplus identity]\label{thm:surplus}
For every $0<a<1$ and $T>0$, the density in \eqref{eq:p} satisfies
$p_a(T)>1$. More precisely,
\begin{align}\label{eq:positive-surplus}
 p_a(T)-1={}&\frac1{\Gamma(b)}
 \int_0^T[k_a(T-t)-k_a(T)]t^{b-2}
          [\phi(t)-\phi(T)]\dd t\notag\\
 &+\frac{k_a(T)}{\Gamma(b)}
 \int_T^\infty t^{b-2}[\phi(T)-\phi(t)]\dd t.
\end{align}
Both integrals are finite and strictly positive.
\end{theorem}
\begin{proof}
For $0<b<1$, the standard subtracted Bose integral is
\begin{equation}\label{eq:zetasubtracted}
 \zeta(b)=\frac1{\Gamma(b)}\int_0^\infty
 t^{b-1}\left(\frac1{e^t-1}-\frac1t\right)\dd t.
\end{equation}
For clarity, its continuation can be justified by splitting at one:
subtract $1/t$ only on $(0,1)$ in the usual zeta integral, add
$1/(b-1)$, and then use
$-\int_1^\infty t^{b-2}\dd t=1/(b-1)$.
The resulting integral converges at both ends for $0<b<1$.

Write $h(t)=1+r(t)$, where $r(t)=(e^t-1)^{-1}-t^{-1}$.
The beta convolution of $k_a$ with $k_b$ is $k_1=1$. Substituting
\eqref{eq:zetasubtracted} into \eqref{eq:p} therefore gives
\begin{align*}
 p_a(T)-1={}&\frac1{\Gamma(b)}
 \int_0^T[k_a(T-t)-k_a(T)]t^{b-1}r(t)\dd t\\
 &-\frac{k_a(T)}{\Gamma(b)}\int_T^\infty t^{b-1}r(t)\dd t.
\end{align*}
All terms are now separately integrable. Since
$t^{b-1}r(t)=t^{b-2}(\phi(t)-1)$, the two contributions from the
constant $-1$ cancel by Lemma~\ref{lem:balance}. Subtracting the
constant $\phi(T)$ against the same balanced signed weight proves
\eqref{eq:positive-surplus}.

On $(0,T)$, $k_a(T-t)>k_a(T)$ because $a<1$, and
$\phi(t)>\phi(T)$. On $(T,\infty)$ the other bracket is also strictly
positive. The local behavior at zero is $O(t^{b-1})$, the singularity
at $T$ is integrable, and the tail has power $t^{b-2}$ with exponent
less than $-1$. This proves finiteness and strict positivity.
\end{proof}

An equivalent form, useful for large $T$, is
\begin{align}\label{eq:Bose-surplus}
 p_a(T)-1={}&\frac1{\Gamma(b)}
 \int_0^T\frac{[k_a(T-t)-k_a(T)]t^{b-1}}{e^t-1}\dd t\notag\\
 &-\frac{k_a(T)}{\Gamma(b)}
 \int_T^\infty\frac{t^{b-1}}{e^t-1}\dd t.
\end{align}
The second term is negative, so this version is not by itself a
positivity proof. Formula~\eqref{eq:positive-surplus}, not a guessed sign
of an asymptotic expansion, supplies that proof.

\subsection{Stronger harmonic and Hurwitz moment inequalities}
\begin{theorem}[Strict critical surplus moments]\label{thm:moments}
For $0<a<1$, put
\begin{equation}\label{eq:Qn}
 Q_n(a)=\frac1a-\frac{H_{n-1}^{(1-a)}}{n^a}-\frac1n,
 \qquad n\ge1.
\end{equation}
Then
\begin{equation}\label{eq:Qmoment}
 Q_n(a)=\int_0^1u^{n-1}[p_a(-\log u)-1]\dd u.
\end{equation}
Every $\Delta^kQ_n(a)$ is strictly positive, including $k=0$.
Every matrix $(Q_{j+k+1})_{j,k=0}^r$ is positive definite, and
$Q_nQ_{n+2}>Q_{n+1}^2$. In particular,
\begin{align}
 \frac{H_{n-1}^{(1-a)}}{n^a}&<\frac1a-\frac1n,
 \label{eq:hbound}\\
 \frac{H_n^{(1-a)}}{(n+1)^a}
 -\frac{H_{n-1}^{(1-a)}}{n^a}&>\frac1{n(n+1)}.
 \label{eq:increment}
\end{align}
\end{theorem}
\begin{proof}
Subtract the Lebesgue moment $1/n$ from the density moment in
Lemma~\ref{lem:measure}. Theorem~\ref{thm:surplus} gives a positive,
nonzero density on every open subinterval of $(0,1)$. Hence
\[
 \Delta^kQ_n=\int_0^1u^{n-1}(1-u)^k[p_a(-\log u)-1]\dd u>0.
\]
For any nonzero polynomial $P$, the integral of $P(u)^2$ against this
density is positive. This is precisely positive definiteness of the
Hankel matrices. Strict Cauchy--Schwarz with the weighted measure
$u^{n-1}[p_a(-\log u)-1]\dd u$ gives strict log-convexity.
The cases $k=0$ and $k=1$ give \eqref{eq:hbound} and
\eqref{eq:increment}. The general Hausdorff-moment mechanism is
classical; see \cite{liu-pego} for the associated generating-function
viewpoint. The new input here is the specific positive surplus density.
\end{proof}

The corresponding continuous identity is
\begin{equation}\label{eq:Qcontinuous}
 Q_a(q)=\frac1a-q^{-a}\bigl[\zeta(1-a)-\zeta(1-a,q)\bigr]-\frac1q
 =\int_0^\infty e^{-qT}[p_a(T)-1]\dd T>0.
\end{equation}
For every $q>0$ and integer $j\ge0$, differentiation gives
$(-1)^jQ_a^{(j)}(q)>0$. Exponential domination at infinity and the
integrable local density justify all derivatives. This strengthens the
prior positivity of the uncorrected critical defect: a full Lebesgue
moment may be subtracted while retaining strict complete monotonicity.

There are also two exact normalization identities:
\begin{equation}\label{eq:surplusnormalizations}
 \int_0^\infty[p_a(T)-1]\dd T=\frac1a,
 \qquad
 \int_0^\infty e^{-T}[p_a(T)-1]\dd T=\frac1a-1.
\end{equation}
For the first, the Hurwitz shift relation and its Taylor expansion at one give
$\zeta(b,q)=q^{-b}+\zeta(b)-b\zeta(b+1)q+O(q^2)$ as $q\downarrow0$.
Equation~\eqref{eq:Qcontinuous} becomes
$Q_a(q)=a^{-1}-b\zeta(b+1)q^b+O(q^{b+1})$.
Monotone convergence against the positive density gives the first identity;
the second is $Q_a(1)=1/a-1$. In particular $a[p_a(T)-1]\dd T$ is a
probability measure on $(0,\infty)$ whose exponential moment
$\E(e^{-T})$ equals $1-a$.

\subsection{A zero-free difference of polylogarithms}
\begin{corollary}[Surplus Stieltjes factorization]\label{cor:factor}
For $0<a<1$ and $z\in\cut$,
\begin{equation}\label{eq:surplusfactor}
 \frac{1-z}{z^2}
 \bigl[F_{a,1-a}(z)-\Li_0(z)+\Li_1(z)\bigr]
 =\int_0^1\frac{(1-u)[p_a(-\log u)-1]}{1-zu}\dd u.
\end{equation}
The right side is positive for real $z<1$ and has strictly positive
imaginary part for $\Impart z>0$. Consequently
$F_{a,1-a}-\Li_0+\Li_1$ has no zero on the slit plane except a double
zero at zero, whose leading coefficient is $2^{-a}-1/2$.
\end{corollary}
\begin{proof}
The endpoint function $\Li_0-\Li_1$ has measure $\delta_1-\dd u$.
Subtract it from \eqref{eq:mu}; the resulting signed measure is
$(1/a-1)\delta_1-[p_a(-\log u)-1]\dd u$ and has total mass zero.
Subtracting $z/(1-z)$ times that zero mass gives
\eqref{eq:surplusfactor}. The numerator defines a strictly positive
measure. The imaginary part of its resolvent is
\[
 (\Impart z)\int_0^1
 \frac{u(1-u)[p_a(-\log u)-1]}{|1-zu|^2}\dd u,
\]
which is strict. The real case is immediate from positive denominators.
Finally, comparing the $z^2$ coefficients yields $2^{-a}-1/2>0$.
\end{proof}

This is a zero-free theorem for a \emph{difference} of functions; it is
not a reannouncement of the earlier zero-free theorem for $F_{a,b}$.
Nor does it decide angular uniqueness below $a+b=1$.

\section{Finite Euler depth: endpoints, lower bounds, and a reversal}
\label{sec:finite}
\begin{proposition}[Endpoint error identities and sharp lower envelope]
\label{prop:RNends}
For each fixed $N\ge1$, $R_N$ extends continuously to $[0,1]$, with
\begin{equation}\label{eq:RNends}
 R_N(0)=J_N^*:=\int_0^1\frac{q_N(u)}{1-u}\dd u,
 \qquad
 R_N(1)=I_N:=\int_0^1q_N(u)\dd u.
\end{equation}
For every $0<a<1$, $R_N(a)>I_N$. Thus $I_N$ is the sharp lower
envelope over the critical line.
\end{proposition}
\begin{proof}
For fixed $T>0$, \eqref{eq:p} and the endpoint beta laws give
\[
 \lim_{a\downarrow0}p_a(T)=\frac1{1-e^{-T}},
 \qquad \lim_{a\uparrow1}p_a(T)=1.
\]
Indeed, $-\zeta(1-a)/\Gamma(a)\to1$ at zero and tends to $1/2$ at
one; the respective beta averages tend to $h(T)$ and $1/2$.
The coefficient $-\zeta(1-a)/\Gamma(a)$ is bounded on the closed
parameter interval after these continuous extensions. Consequently,
with an absolute constant,
\begin{equation}\label{eq:pglobalbound}
 0<p_a(T)\le C(1+T^{-1})\quad(T>0,\ 0<a<1).
\end{equation}
For $T=-\log u$, the right side is at most $C'(1-u)^{-1}$.
The weight $q_N(u)/(1-u)$ is integrable for $N\ge1$, so dominated
convergence proves the endpoint limits and continuity. Strict surplus
positivity proves $R_N(a)>I_N$.
\end{proof}

The limiting density at $a=0$ has infinite total mass. Only its filtered
integral in \eqref{eq:RNends} is finite; it must not be substituted as a
finite signed measure into \eqref{eq:mu}.

There are elementary finite recurrences for these endpoint values. Let
\begin{equation}\label{eq:BNpoly}
 B_j^{\rm pol}=\int_0^1(1-u^2)^j\dd u
 =\frac{4^j(j!)^2}{(2j+1)!}.
\end{equation}
Then
\begin{align}
 I_0&=\frac\pi4,& I_{N+1}&=2I_N-B_N^{\rm pol}\quad(N\ge0),
 \label{eq:Irec}\\
 J_1^*&=\frac\pi4+\frac{\log2}{2},&
 J_{N+1}^*&=2J_N^*-B_{N-1}^{\rm pol}-\frac1{2N}\quad(N\ge1).
 \label{eq:Jrec}
\end{align}
To prove these, use $q_{N+1}=2q_N-(1-u^2)^N$ and
$(1-u^2)^N/(1-u)=(1+u)(1-u^2)^{N-1}$.
The recurrences are exact identities, but their subtractive form may be
poor for floating-point evaluation at large $N$.

\begin{theorem}[A certified finite-depth reversal]\label{thm:reversal}
Strictly,
\begin{align}
 0.084208864068217362453965253610
 &<R_{128}(1/4)
 <0.084208864068217362453965253612,
 \label{eq:certquarter}\\
 0.084523512538326867477886761599
 &<R_{128}(1/2)
 <0.084523512538326867477886761601.
 \label{eq:certhalf}
\end{align}
In particular $R_{128}(1/2)>R_{128}(1/4)$, whereas
$R_1(1/2)<R_1(1/4)$.
\end{theorem}
\begin{proof}
The last inequality is also an instance of Theorem~\ref{thm:sharp}.
The two displayed enclosures follow from the standard-library-only
certificate \texttt{code/certify.py}. Its finite arithmetic and tail
contract are described in Section~\ref{sec:verification}; a verifier
replays integer-root acceptance inequalities and outward interval
operations. In particular the assertion is not inferred from decimal
agreement between two floating-point calculations.
\end{proof}

This counterexample refutes a \emph{stronger extrapolation}, namely that
$R_N(a)$ decreases in $a$ for each $N$. The incoming report does not
assert that extrapolation. Its actual conjecture concerns $R_1$, or
equivalently the optimal constant after taking the supremum over $N$;
that conjecture is proved, not contradicted, here.

\section{Large-logarithm structure of the critical density}
\label{sec:densityasymp}
Put $b=1-a$ and define the positive coefficients
\begin{equation}\label{eq:alphaj}
 \alpha_j(b)=\frac{(b)_j^2\zeta(b+j)}{j!\Gamma(1-b)},
 \qquad j\ge1,
 \qquad A(b)=\alpha_1(b)=\frac{b^2\zeta(1+b)}{\Gamma(1-b)}.
\end{equation}
Here $(b)_j=b(b+1)\cdots(b+j-1)$ is a rising factorial.

\begin{theorem}[All-orders density expansion]\label{thm:densityall}
For fixed $0<b<1$ and every integer $M\ge1$,
\begin{equation}\label{eq:densityall}
 p_{1-b}(T)=1+\sum_{j=1}^M\alpha_j(b)T^{-b-j}
             +O(T^{-b-M-1})\quad(T\to\infty).
\end{equation}
The remainder is uniform when $b$ ranges over a compact subinterval of
$(0,1)$. This is a Poincar\'e asymptotic expansion, not a claim of
convergence for a fixed $T$ or a globally signed truncation theorem.
\end{theorem}
\begin{proof}
Write \eqref{eq:Bose-surplus} as
\begin{align}\label{eq:scaledBose}
 p_{1-b}(T)-1=\frac{T^{-b}}{\Gamma(1-b)\Gamma(b)}
 \left\{\int_0^T\frac{[(1-t/T)^{-b}-1]t^{b-1}}{e^t-1}\dd t
 -\int_T^\infty\frac{t^{b-1}}{e^t-1}\dd t\right\}.
\end{align}
On $0\le t\le T/2$, expand
$(1-t/T)^{-b}-1=\sum_{j=1}^M(b)_j(t/T)^j/j!
+O((t/T)^{M+1})$. The remainder constant is uniform on compact
$b$-sets. Integrating termwise and extending each exponentially decaying
moment to infinity introduces exponentially small errors. The moments
are $\Gamma(b+j)\zeta(b+j)$, so their coefficients are exactly
\eqref{eq:alphaj}. On $T/2<t<T$, the integrable singularity
$(T-t)^{-b}$ is multiplied by an exponentially small Bose factor;
its contribution and the last integral are smaller than every displayed
power, uniformly on the same compact sets. This proves the result.
\end{proof}

The maximizer will approach a parameter endpoint. A merely compact-uniform
statement is insufficient for optimizing over all $a$. The following
estimate supplies the missing endpoint uniformity.

\begin{lemma}[A uniform first-term estimate]\label{lem:uniformdensity}
There are absolute constants $C,T_0$ such that for all $0<b<1$ and
$T\ge T_0$,
\begin{equation}\label{eq:uniformdensity}
 \left|p_{1-b}(T)-1-A(b)T^{-1-b}\right|
 \le C A(b)T^{-2-b}+C e^{-T/2}.
\end{equation}
One may take $T_0=2$ with a sufficiently large absolute $C$.
\end{lemma}
\begin{proof}
For $0\le s\le1/2$ and $0<b<1$, Taylor's formula gives
\[
 0\le(1-s)^{-b}-1-bs\le8bs^2.
\]
Use this in \eqref{eq:scaledBose} on $(0,T/2)$.
Its leading integral, when extended to infinity, is
$A(b)T^{-1-b}$. The remainder is at most
\[
 \frac{8b\Gamma(b+2)\zeta(b+2)}
 {\Gamma(1-b)\Gamma(b)}T^{-b-2}
 =8(b+1)\frac{\zeta(b+2)}{\zeta(b+1)}A(b)T^{-b-2}
 \le16A(b)T^{-b-2}.
\]
The omitted tail of the leading moment is $O(e^{-T/2})$ uniformly:
for $x\ge1$ and $0<b<1$,
$\int_x^\infty t^be^{-t}\dd t\le2x^be^{-x}$, while reciprocal
gamma factors are bounded on the parameter interval.

For the singular contribution on $T/2<t<T$, use
$(e^t-1)^{-1}\le(1-e^{-1})^{-1}e^{-t}$ and retain the complete beta
weight rather than bounding its endpoint exponent separately. The
term containing $(1-t/T)^{-b}$ is bounded by a constant times
\[
 \frac{e^{-T/2}}{\Gamma(1-b)\Gamma(b)}
 \int_0^T(T-t)^{-b}t^{b-1}\dd t=e^{-T/2}.
\]
The remaining nonsingular tails in \eqref{eq:scaledBose} are bounded
by the same type of estimate, using $T\ge2$. This avoids a spurious
factor diverging as $b\uparrow1$. Combining the bounds proves
\eqref{eq:uniformdensity}.
\end{proof}

\section{Uniform Euler asymptotics and the drifting extremizer}
\label{sec:optimization}
For the remainder of this section set
\begin{equation}\label{eq:Lell}
 L=\tfrac12\log N,\qquad \ell=\log L,\qquad
 B_N=\frac{\sqrt\pi}{2\sqrt N}.
\end{equation}
Statements using $\ell$ concern $N\to\infty$; they are not proposed as
accurate low-depth formulas.

\subsection{A uniform localization lemma}
\begin{lemma}[Gaussian logarithmic averaging]\label{lem:logaverage}
Uniformly for $1\le s\le3$,
\begin{equation}\label{eq:logaverage}
 \int_0^{N^{-1/4}}q_N(u)(-\log u)^{-s}\dd u
 =B_NL^{-s}\bigl(1+O(L^{-1})\bigr).
\end{equation}
Also
\begin{equation}\label{eq:INasymp}
 I_N=B_N\left(1-\frac7{8N}+O(N^{-2})\right).
\end{equation}
\end{lemma}
\begin{proof}
Set $u=t/\sqrt N$. Then $-\log u=L-\log t$ and
$0<t<N^{1/4}$. On this interval
\[
 q_N(t/\sqrt N)=e^{-t^2}
 +O\left(N^{-1}(t^2+t^4)e^{-t^2}\right),
\]
in the sense of an absolute pointwise bound with a fixed constant for
large $N$. This follows from
$-v-v^2/[2(1-v)]\le\log(1-v)\le-v$ and
$1-e^{-x}\le x$, with $v=t^2/N\le N^{-1/2}$.
On $e^{-L/2}\le t\le e^{L/2}$, the mean-value theorem gives, uniformly
for $1\le s\le3$,
\[
 \left|\left(1-\frac{\log t}{L}\right)^{-s}-1\right|
 \le C\frac{|\log t|}{L}.
\]
The Gaussian logarithmic moment is integrable. On $0<t<e^{-L/2}$
the normalized factor lies between zero and one and the interval has
length $e^{-L/2}$. Gaussian tails above $e^{L/2}$ are negligible.
This proves \eqref{eq:logaverage} with constants independent of $s$.

For \eqref{eq:INasymp}, expand one term further on a central interval
and bound its complement by Gaussian decay:
\[
 q_N(t/\sqrt N)=e^{-t^2}
 \left[1-\frac{t^2+t^4/2}{N}+O(N^{-2}P(t))\right],
\]
where a fixed polynomial $P$ suffices on a growing central interval.
The normalized half-Gaussian moments of $t^2$ and $t^4$ are $1/2$ and
$3/4$. Hence the coefficient is $1/2+(3/4)/2=7/8$.
Taylor remainders on, for example, $t\le N^{1/12}$ and exponential
bounds on the complement justify the integrated $O(N^{-2})$ term.
\end{proof}

\begin{theorem}[Uniform critical Euler expansion]\label{thm:uniformEuler}
Uniformly for every $0<b<1$ as $N\to\infty$,
\begin{equation}\label{eq:uniformEuler}
 R_N(1-b)=I_N+B_NA(b)L^{-1-b}\bigl(1+O(L^{-1})\bigr)
                   +O(N^{-3/4}).
\end{equation}
More explicitly, the absolute difference after subtracting
$I_N+B_NA(b)L^{-1-b}$ is at most
$C B_NA(b)L^{-2-b}+C N^{-3/4}$ with an absolute constant.
In particular,
\begin{equation}\label{eq:uniformlimit}
 \sup_{0<a<1}\left|\frac{R_N(a)}{B_N}-1\right|\longrightarrow0.
\end{equation}
\end{theorem}
\begin{proof}
By \eqref{eq:Rintegral}, the difference $R_N-I_N$ is the integral of
$q_N(u)[p_a(-\log u)-1]$. First restrict to $u\le N^{-1/4}$, so that
$T=-\log u\ge L/2\ge2$ for large $N$. Apply
Lemma~\ref{lem:uniformdensity}. The leading term is evaluated by
Lemma~\ref{lem:logaverage} with $s=1+b$; the algebraic error uses
$s=2+b$. The exponential density error integrates to at most
\[
 C\int_0^1q_N(u)u^{1/2}\dd u
 \le C\int_0^\infty e^{-Nu^2}u^{1/2}\dd u=O(N^{-3/4}).
\]
For the complementary interval, use \eqref{eq:pglobalbound} and
\[
 \frac{q_N(u)}{1-u}
 =\frac{1+u}{1+u^2}(1-u^2)^{N-1}
 \le2e^{-(N-1)N^{-1/2}}\qquad(u\ge N^{-1/4}).
\]
Its total contribution is exponentially small uniformly in $a$.
This proves the explicit remainder. Finally, $A(b)$ is bounded on
$(0,1)$ and extends continuously with value zero at both endpoints,
so \eqref{eq:uniformlimit} follows from \eqref{eq:INasymp}.
\end{proof}

\subsection{Fixed-parameter coefficients to all orders}
Let
\begin{equation}\label{eq:logmoments}
 m_r=\frac{2}{\sqrt\pi}\int_0^\infty e^{-t^2}(\log t)^r\dd t
 =\frac{\Gamma^{(r)}(1/2)}{2^r\Gamma(1/2)}.
\end{equation}
In particular $m_0=1$, $m_1=-\gamma/2-\log2$, and
$m_2=m_1^2+\pi^2/8$.
For $m\ge1$ define the explicit zeta--gamma coefficient
\begin{equation}\label{eq:Dm}
 D_m(b)=\sum_{j=1}^m\alpha_j(b)
 \frac{(b+j)_{m-j}}{(m-j)!}\,m_{m-j}.
\end{equation}

\begin{corollary}[All logarithmic orders]\label{cor:allEuler}
For fixed $0<b<1$ and any integer $M\ge1$,
\begin{equation}\label{eq:allEuler}
 R_N(1-b)=I_N+B_N\sum_{m=1}^M D_m(b)L^{-b-m}
                     +O\bigl(B_NL^{-b-M-1}\bigr).
\end{equation}
The result is uniform on compact $b$-subintervals of $(0,1)$.
The first three coefficients are
\begin{align*}
 D_1&=\alpha_1,\\
 D_2&=\alpha_2+(b+1)m_1\alpha_1,\\
 D_3&=\alpha_3+(b+2)m_1\alpha_2
                +\tfrac12(b+1)(b+2)m_2\alpha_1.
\end{align*}
\end{corollary}
\begin{proof}
Insert Theorem~\ref{thm:densityall} into the localized error integral.
For each density term expand
\[
 (L-\log t)^{-b-j}
 =L^{-b-j}\sum_{r=0}^{M-j}
 \frac{(b+j)_r}{r!}\frac{(\log t)^r}{L^r}
 +\text{remainder}.
\]
On $|\log t|\le L/2$ Taylor's remainder is controlled by the next
Gaussian logarithmic moment. The omitted small-$t$ interval has
exponentially small length, and the large-$t$ interval has Gaussian
decay. The $q_N$ replacement error is polynomially small in $N$, hence
smaller than every fixed inverse logarithmic order. Integrate termwise
using \eqref{eq:logmoments}, and collect terms with $j+r=m$ to obtain
\eqref{eq:Dm}. Compact-uniform constants follow from the same estimates.
\end{proof}

This finite coefficient formula gives an entire family of proved
asymptotic identities. The positivity of the density coefficients
$\alpha_j$ does not imply that every $D_m$ is positive: logarithmic
Gaussian moments can have either sign.

\subsection{Maximization over the critical line}
Let $R_N(0)$ and $R_N(1)$ have their continuous endpoint meanings from
Proposition~\ref{prop:RNends}, and define
\begin{equation}\label{eq:CN}
 C_N=\max_{0\le a\le1}R_N(a).
\end{equation}
A maximizer exists by continuity. Its value is also the supremum over
the open critical line.

\begin{theorem}[Drift of every worst-case parameter]\label{thm:drift}
For all sufficiently large $N$, every maximizer $a_N$ in
\eqref{eq:CN} lies in $(0,1)$. As $N\to\infty$,
\begin{align}
 (1-a_N)\ell&\longrightarrow1,
 \label{eq:drift}\\
 C_N&=B_N\left[1+\frac{1+o(1)}{eL\ell}\right].
 \label{eq:CNasymp}
\end{align}
No uniqueness of the finite-$N$ maximizer is assumed or concluded.
\end{theorem}
\begin{proof}
Extend $A$ to $[0,1]$ by zero at both endpoints. Since the zeta pole
has residue one and $\Gamma(1)=1$,
\begin{equation}\label{eq:Afirst}
 A(b)=b(1+O(b))\quad(b\downarrow0),
 \qquad A(b)\le Cb\quad(0\le b\le1).
\end{equation}
Define $f_\ell(b)=A(b)e^{-\ell b}$. We claim
\begin{equation}\label{eq:surrogatemax}
 \max_{0\le b\le1}f_\ell(b)\sim\frac1{e\ell},
 \qquad \ell b_\ell\to1
\end{equation}
for every maximizing $b_\ell$ and also every asymptotically maximizing
sequence. Testing $b=1/\ell$ gives the lower estimate. For
$t=\ell b$ in a fixed compact interval,
$\ell f_\ell(t/\ell)\to te^{-t}$ uniformly. The upper bound in
\eqref{eq:Afirst} makes the contribution for large $t$ uniformly small,
and the same bound makes the contribution for small $t$ small.
The function $te^{-t}$ has a unique maximum $1/e$ at $t=1$, proving
\eqref{eq:surrogatemax}.

Theorem~\ref{thm:uniformEuler}, after multiplication by $L/B_N$, is
\begin{equation}\label{eq:normalizedopt}
 \frac{L}{B_N}\bigl(R_N(1-b)-I_N\bigr)
 =f_\ell(b)+O\bigl(f_\ell(b)/L\bigr)+O(LN^{-1/4}),
\end{equation}
uniformly on the open interval. At $a=1$ the left side is zero.
At $a=0$, a direct substitution $u=t/\sqrt N$ gives
\begin{equation}\label{eq:JminusI}
 J_N^*-I_N
 =\int_0^1q_N(u)\frac{u}{1-u}\dd u
 =\frac1{2N}+O(N^{-3/2}).
\end{equation}
Thus \eqref{eq:normalizedopt} is valid at the endpoints as well with
an error of the stated or smaller size. Its errors are $o(1/\ell)$,
whereas the largest surrogate value is asymptotic to $1/(e\ell)$.
Every true maximizer is therefore asymptotically maximizing for the
surrogate and obeys \eqref{eq:drift}. It cannot be an endpoint for
large $N$. The maximum excess is
$B_N/(eL\ell)(1+o(1))$. Finally
$I_N/B_N=1+O(N^{-1})$ proves \eqref{eq:CNasymp}.
\end{proof}

A useful consequence of the same proof is that, for \emph{every fixed}
$0<a<1$, $R_N(a)>R_N(0)$ once $N$ is sufficiently large (depending on
$a$). Indeed its positive excess over $I_N$ has order
$N^{-1/2}L^{-1-b}$, while the $a=0$ excess has order $N^{-1}$.
This supplies an analytic explanation for the finite-depth reversal,
independent of the particular certified pair.

\subsection{A Stieltjes-constant correction}
Use the convention
\[
 \zeta(1+b)=\frac1b+\gamma-\gamma_1b+O(b^2)
\]
for the first Stieltjes constant $\gamma_1$ \cite{dlmf-zeta}.
The gamma expansion \cite{dlmf-gamma} gives
\begin{equation}\label{eq:Acubic}
 A(b)=b-db^3+O(b^4),\qquad
 d=\gamma_1+\frac{\gamma^2+\zeta(2)}2
 =0.916240149844295830534809275626\ldots.
\end{equation}
In particular the quadratic coefficient vanishes exactly.

\begin{theorem}[Refined drift and optimal excess]\label{thm:refined}
For every maximizer $a_N$ and the notation \eqref{eq:Lell},
\begin{align}
 1-a_N&=\frac1\ell-\frac{2d}{\ell^3}+O(\ell^{-4}),
 \label{eq:refinedlocation}\\
 C_N&=B_N\left[1+\frac1{eL\ell}
       \left\{1-\frac d{\ell^2}+O(\ell^{-3})\right\}\right].
 \label{eq:refinedvalue}
\end{align}
The location estimate applies to all maximizing choices, even if some
finite depth has more than one maximizer.
\end{theorem}
\begin{proof}
Expansion of the reciprocal gamma factor in \eqref{eq:alphaj} gives
\[
 \Gamma(1-b)^{-1}
 =1-\gamma b+\tfrac12(\gamma^2-\zeta(2))b^2+O(b^3),
\]
which proves \eqref{eq:Acubic}. Near $b=0$ it follows that
\[
 \frac{A'(b)}{A(b)}=\frac1b-2db+O(b^2).
\]
By \eqref{eq:surrogatemax}, all surrogate maximizers eventually lie
in $3/(4\ell)<b<5/(4\ell)$. On that interval $f_\ell''(b)$ is
negative and bounded above by $-c\ell$ for some $c>0$: differentiating
$A(b)e^{-\ell b}$ and using $A(b)=b+O(b^3)$ gives
$f_\ell''(b)=e^{-\ell b}(\ell^2b-2\ell+O(\ell^{-1}))$.
Hence there is one surrogate maximum $b_\ell^*$ in this interval.
Its equation $A'/A=\ell$ implies
\[
 b_\ell^*=\ell^{-1}-2d\ell^{-3}+O(\ell^{-4}),
 \qquad
 f_\ell(b_\ell^*)=\frac1{e\ell}
             \bigl(1-d\ell^{-2}+O(\ell^{-3})\bigr).
\]

It remains to transfer location, not just value, to the true error.
Theorem~\ref{thm:drift} places every true maximizer in the same
interval for large $N$. On it the uniform error in
\eqref{eq:normalizedopt} is
$O((\ell L)^{-1}+LN^{-1/4})$.
Comparing the true maximum with the value at $b_\ell^*$ bounds the
surrogate value loss by twice this error. Strict concavity therefore
implies
\[
 |(1-a_N)-b_\ell^*|^2
 \le C\left(\frac1{\ell^2L}+\frac{LN^{-1/4}}\ell\right).
\]
Since $L=e^\ell$ and $N=e^{2L}$, this location difference is smaller
than every fixed inverse power of $\ell$. This proves
\eqref{eq:refinedlocation} without differentiating the true error
remainder. The same uniform comparison transfers the surrogate maximum
value with an error smaller than any displayed inverse-$\ell$ order;
\eqref{eq:INasymp} then gives \eqref{eq:refinedvalue}.
\end{proof}

\begin{remark}[Why the two extremal statements are compatible]
The sharp constant after optimizing over \emph{both} $a$ and $N$ is
$\Cstar$, approached at $a\downarrow0,N=1$.
At a large \emph{prescribed} depth $N$, the maximum instead occurs at
an interior parameter approaching $a=1$. The rate is inverse
$\log\log N$, more precisely \eqref{eq:refinedlocation} with
$\ell=\log(\frac12\log N)$. This slow rate makes the limiting formulas
poor low-depth predictors. The statements concern different
optimizations; neither reverses or weakens Theorem~\ref{thm:sharp}.
\end{remark}

\section{Exact certificates and independent computations}
\label{sec:verification}
\subsection{Arithmetic contract of the rational verifier}
All certified quantities use the already proved bound
$0<E_K-g<(1/a)2^{-K}$ with $K=256$. In particular, the verifier does
not assume the new constant $\Cstar$ to certify its examples.
The finite sum is formed with the identity
\begin{equation}\label{eq:Eulerweights}
 E_N=2^{-N}\sum_{n=0}^{N-1}(-1)^nf(n)
                   \sum_{j=n+1}^N\binom Nj.
\end{equation}
To prove it, substitute the finite definition of $d_k$ and use
\[
 \sum_{k=n}^{N-1}2^{N-k-1}\binom kn
 =\sum_{j=n+1}^N\binom Nj.
\]
Pascal's identity proves the coefficient relation by induction.
This avoids constructing a large interval difference table.

For an inverse rational power $n^{-p/q}$, choose $S=10^{210}$ and
compute an integer proposal $r$ for the $q$th root of $n^pS^q$.
The verifier accepts it only after checking
\begin{equation}\label{eq:rootcertificate}
 r^q\le n^pS^q<(r+1)^q.
\end{equation}
If the left equality does not hold, $S/(r+1)<n^{-p/q}<S/r$.
In either case these rational endpoints are rounded outward to the
grid $S^{-1}\mathbb Z$ using integer division. Positive products,
signed linear combinations, and division by positive integers are
also rounded outward. The root-finding iteration is a proposal method;
\eqref{eq:rootcertificate}, not the iteration history, is its guarantee.

An enclosure for $g$ is obtained by subtracting $(1/a)2^{-256}$ from
the lower endpoint of the finite $E_{256}$ enclosure and leaving the
upper endpoint unchanged. Subtract this $g$ interval from the finite
$E_N$ interval and multiply by $2^N$ to obtain $R_N$. At $N=128$ the
analytic tail width is at most $(1/a)2^{-128}$, well below the displayed
30-place comparison. Finally the program tests the human-readable
rational endpoints in \eqref{eq:certquarter}--\eqref{eq:certhalf}
directly by integer cross multiplication.

\subsection{Recorded exact checks}
The supplied run passes 2,555 integer-root acceptance checks. It
certifies Gaussian values and six scaled-error depths for each of
$a=1/10,1/4,1/2,3/4,9/10$, hence 30 positive-error enclosures. It also
checks 56 strict surplus finite differences for each parameter, for
280 such tests in total. These finite tests illustrate
Theorem~\ref{thm:moments}; the theorem's all-index content comes from
the positive integral proof, not from extrapolating the test grid.
The exact difference certificate is
\begin{equation}\label{eq:reversaldiff}
 \begin{split}
 0.000314648470109505023921507988
 &<R_{128}(1/2)-R_{128}(1/4)\\
 &<0.000314648470109505023921507989.
 \end{split}
\end{equation}
The JSON stores outward intervals rather than unsupported point values.
A separate SymPy script checks eight algebraic identities and 528 finite
instances of the binomial weight identity. It does not check analytic
continuation, dominated convergence, or the uniformity of asymptotic
remainders; those arguments are in the article.

\subsection{Independent numerical diagnostics}
The floating-point program uses 180 decimal digits for finite Euler
calculations. Five Gaussian values are separately checked by
beta--gamma quadrature, which evaluates \eqref{eq:Gtransport} rather
than an Euler difference table. Nine density tests compare the beta
average \eqref{eq:p} with the regularized Bose expression
\eqref{eq:Bose-surplus} at $a=1/4,1/2,3/4$ and $T=2,10,50$.
These checks also compare one- and two-term density asymptotics.
They are diagnostics, not rigorous quadrature error bounds.

A grid scan followed by bounded local optimization gives the following
\emph{proposals}; uniqueness and global optimality at these particular
depths are not certified by the optimizer.
\begin{center}
\small
\begin{tabular}{@{}rrr@{}}
\toprule
$N$ & proposed maximizing $a$ & proposed $\max R_N$\\
\midrule
8   & 0.10599483 & 0.359993461628\\
16  & 0.21532658 & 0.249985265562\\
32  & 0.29829873 & 0.173968759177\\
64  & 0.36226668 & 0.121253639850\\
128 & 0.41230987 & 0.084666552368\\
256 & 0.45206073 & 0.059234316775\\
\bottomrule
\end{tabular}
\end{center}
The movement is consistent with the proved drift, but the proof is
independent of the table. The program additionally optimizes the
surrogate $A(b)e^{-\ell b}$ for larger $\ell$; those surrogate results
are explicitly labelled and are not presented as astronomical-$N$
computations of the original error.

\subsection{Reproduction}
From the package directory, run
\begin{verbatim}
python code/certify.py
python code/symbolic_checks.py
python code/diagnostics.py
python code/build_article.py
\end{verbatim}
Only the first command is needed to replay the rational certificates;
it uses the Python standard library. The optional diagnostics require
mpmath, NumPy, SciPy, and SymPy at the versions listed in
\texttt{requirements-diagnostics.txt}. The build requires a local
\texttt{pdflatex}. All scripts resolve their paths relative to their
own location, accept no external numerical oracle, and write their
results to the package's \texttt{data/} directory. Run them on a copy
when preserving a delivered evidence snapshot.

\section{Audit conclusions and integration}
\label{sec:audit}
\subsection{What should change in the claim ledger}
The incoming report's Conjecture~12.3 can be promoted to a theorem with
Theorem~\ref{thm:sharp} as its proof reference. Its question about sharp
$N$-dependence is answered by Theorems~\ref{thm:uniformEuler},
\ref{thm:drift}, and \ref{thm:refined}. The new proof route relies on
the existing positive double kernel and critical measure, not on any
unproved numerical relation among special values.

No error in the \emph{stated} critical measure or Euler bound of the
reviewed incoming report was established. The relevant corrections are
to possible extensions or editorial wording: do not replace the
critical measure by a density without its atom; do not describe its
boundary series as ordinarily convergent; do not omit the factor $X$
from the new rate-one beta--gamma coupling; and do not extend the
proved decrease of $R_1$ to the functions $R_N$ individually.
The last overextension is now excluded by an exact certificate.

The asymptotic expansion \eqref{eq:densityall} must not be advertised as
a convergent global series or as a signed remainder theorem. For
parameter optimization use \eqref{eq:uniformdensity}, which is uniform
up to both parameter endpoints, rather than the compact-uniform
fixed-order expansion. Finally, the stationary point of the surrogate
is not the exact finite-$N$ optimizer; Theorem~\ref{thm:refined} proves
their asymptotic proximity and the numerical records keep them separate.

\subsection{Suggested repository placement}
The complete delivery can be placed under
\begin{center}\small
\nolinkurl{Analysis/Polylogarithms/docs/reports/critical-euler-transport/}.
\end{center}
Its thematic predecessor is the incoming real-order-threshold report.
The manuscript insertion supplied in \texttt{integration/} uses the
prefix \texttt{criticaleuler:} and is a summary with references to the
full proofs, not a replacement for them. The main proof of the sharp
constant belongs near the signed-kernel and certified-evaluation
material; the positive surplus moments connect naturally to the
Hurwitz-jet chapters. Intake should preserve the original delivery and
record the change of conjecture status in the programme ledger.

The existing $S_6$ identity, subcritical angular uniqueness, and the
broad normalized-radius conjecture are not resolved by this article.
No new independence, minimal-depth, or transcendence assertion follows
from these positive kernels. This distinction is especially important
because a new numerical evaluation method does not by itself reduce a
constant to a prescribed arithmetic period basis.

\section{Further research questions}
\label{sec:future}
\subsection{Uniqueness and monotonic motion of the finite-depth maximizer}
The asymptotic theory localizes every maximizer but does not exclude
several extremely close maxima at a finite depth. The numerical samples
suggest a stronger statement.
\begin{conjecture}[Eventual unimodality]\label{conj:unimodal}
For every sufficiently large integer $N$, $a\mapsto R_N(a)$ has exactly
one critical point in $(0,1)$, a strict global maximum. The maximizing
outer order increases strictly with $N$ once $N$ is sufficiently large.
\end{conjecture}
A route to proof is to differentiate the balanced density identity in
the shape parameter and establish uniform derivative remainders. The
function-value argument in Theorem~\ref{thm:refined} deliberately avoids
such estimates and therefore cannot supply uniqueness. A useful finite
problem is to determine the first depth at which the maximum leaves
$a=0$, with certified endpoint derivative signs and an exhaustive
zero count for the derivative.

\subsection{An explicit nonasymptotic optimal-depth bound}
Theorem~\ref{thm:uniformEuler} gives absolute uniform constants but does
not optimize or enumerate them. An implementable bound of the form
\[
 R_N(a)\le\frac{\sqrt\pi}{2\sqrt N}
 \left(1+\frac{K}{\log N\,\log\log N}\right)
\]
with explicit $K$ and a verified threshold would improve practical
certificates for all critical orders at once. Below that threshold one
could combine the sharp all-depth constant with a finite exact
parameter enclosure. This requires interval control over a continuum
of $a$, not merely a table of rational sample orders.

\subsection{The exact radial domain of order transport}
Determine for which pairs $(w,\rho)$ the strict decrease of
$G_{a,w-a}(\rho)$ holds for all $a\in(0,w)$. Our theorem covers
$0<\rho\le\sqrt3$ for every $w>0$. Beyond that radius the pointwise
kernel changes derivative sign, but the beta--gamma average may preserve
the comparison in a larger domain. A proof or a rigorously enclosed
counterexample should distinguish these two mechanisms.

\subsection{Other Gaussian filters and cyclotomic points}
The filter $(1-u^2)^N/(1+u^2)$ is particularly suited to the Gaussian
point. Root-of-unity filters of higher order lead to different localized
moments and may change the worst-order drift. It is natural to ask
whether balanced critical densities give similarly sharp constants and
explicit Stieltjes-constant corrections for Eisenstein or other
cyclotomic evaluations. The positivity of the present filter cannot
simply be presumed for a complex-valued generalization.

\subsection{Higher-depth critical defects}
Replacing $H_{n-1}^{(b)}$ by nested harmonic sums should introduce more
than one critical balance. A precise objective is to identify the
largest elementary positive moment measure that can be subtracted from
the corresponding defect while retaining positivity. At depth two the
answer supplied here includes the entire Lebesgue measure. At higher
depth, an endpoint atom and an unproved sign pattern must not be
confused with a fully established signed Hausdorff representation.

\subsection{Remaining arithmetic and geometric questions}
The $S_6$ reduction and the subcritical angular-zero conjecture remain
separate worthwhile targets. The present surplus factorization may
help compare critical functions, but it does not construct a finite
signed measure below the known threshold. On the computational side,
the finite integer-root verifier is a small candidate for formal
verification; the analytic layer would additionally require beta--gamma
integration, signed measure transforms, and uniform asymptotics.

\section*{Conclusion}
\addcontentsline{toc}{section}{Conclusion}
The critical constant is now determined by a strict beta-order comparison,
not by numerical optimization. A balanced Bose-kernel identity reveals a
positive density surplus, giving stronger harmonic moment inequalities
and a new zero-free difference factorization. That same surplus explains
why the worst finite-depth Euler parameter moves in the opposite
direction from the endpoint controlling the sharp all-depth constant.
The uniform error theorem and its Stieltjes-constant refinement identify
this movement quantitatively, while a rational certificate exhibits the
reversal at a concrete finite depth. These results are a continuation of
the inspected real-order programme, with unresolved arithmetic and
geometric questions left explicitly separate.

\appendix
\section{Claim-status summary}
\small
\begin{longtable}{@{}>{\raggedright\arraybackslash}p{.38\textwidth}>{\raggedright\arraybackslash}p{.56\textwidth}@{}}
\toprule
Claim & Evidence and status\\\midrule\endhead
Double kernel; critical measure; Euler remainder & Prior analytic identities, rederived where needed and attributed.\\[3pt]
Fixed-total-order transport and comparisons & Proved in Section~\ref{sec:transport}; imaginary radii through $\sqrt3$ only.\\[3pt]
Incoming Conjecture 12.3 & Proved, including strict monotonicity and sharp uniform constant.\\[3pt]
Strict surplus density and positive moments & Proved by a balanced integral, with all-index and continuous-parameter consequences.\\[3pt]
Zero-free surplus difference & Proved for the principal branch on the slit plane.\\[3pt]
Uniform Euler expansion and extremizer drift & Proved with endpoint-uniform error; no uniqueness premise.\\[3pt]
Refined Stieltjes-constant coefficient & Proved by surrogate comparison and curvature, not a fitted numerical coefficient.\\[3pt]
Finite-depth reversal & Exact outward rational certificate using a proved infinite-tail bound.\\[3pt]
Quadrature and numerical optimizer tables & Diagnostics only; no interval-certified global optimum claim.\\[3pt]
Eventual unimodality and monotone maximizing parameter & Conjectural.\\[3pt]
$S_6$, subcritical angular uniqueness, broad normalized-radius claim & Not settled here.\\[3pt]
Global novelty; proof-assistant formalization; exhaustive source audit & Not claimed.\\
\bottomrule
\end{longtable}
\normalsize

\addcontentsline{toc}{section}{References}
\begin{thebibliography}{9}
\bibitem{manuscript}
ProveIt Contributors, \emph{Polylogarithms and their Arithmetic Bridges:
Special Values, Zeta Jets, Lattices and Experimental Discovery}.
Canonical manuscript, repository \texttt{VladimirReshetnikov/ProveIt},
revision \basecommit.
\url{https://github.com/VladimirReshetnikov/ProveIt/tree/a0a90ef31877f98be437191c48b46f02d5456867/Analysis/Polylogarithms/docs/manuscript}.

\bibitem{threshold}
ProveIt Contributors, \emph{The Sharp Real-Order Threshold for Harmonic
Polylogarithms: Endpoint atoms, zero geometry, and certified fractional-order
identities}, incoming research delivery, 9 October 2026.
\texttt{ProveIt\_RealOrder\_Threshold.zip}, especially Theorem~3.1,
Section~9, Proposition~12.2, and Conjecture~12.3.
\url{https://github.com/VladimirReshetnikov/ProveIt/blob/a0a90ef31877f98be437191c48b46f02d5456867/docs/incoming/ProveIt_RealOrder_Threshold.zip}.

\bibitem{dlmf-beta}
NIST Digital Library of Mathematical Functions, \S5.12,
\emph{Beta Function}, especially Euler's beta integral 5.12.1.
Version 1.2.8; accessed 9 October 2026 (Pacific time).
\url{https://dlmf.nist.gov/5.12}.

\bibitem{dlmf-hurwitz}
NIST Digital Library of Mathematical Functions, \S25.11,
\emph{Hurwitz Zeta Function}, especially its regularized integral
representations. \url{https://dlmf.nist.gov/25.11}.

\bibitem{dlmf-polylog}
NIST Digital Library of Mathematical Functions, \S25.12,
\emph{Polylogarithms}. \url{https://dlmf.nist.gov/25.12}.

\bibitem{dlmf-zeta}
NIST Digital Library of Mathematical Functions, \S25.2,
\emph{Definition and Expansions}, including the Laurent expansion and
Stieltjes constants. \url{https://dlmf.nist.gov/25.2}.

\bibitem{dlmf-gamma}
NIST Digital Library of Mathematical Functions, \S5.7,
\emph{Series Expansions} of the gamma function.
\url{https://dlmf.nist.gov/5.7}.

\bibitem{liu-pego}
J.-G. Liu and R. L. Pego, \emph{On generating functions of Hausdorff
moment sequences}, Transactions of the American Mathematical Society
\textbf{368} (2016), 8499--8518. Consult the updated arXiv version and
its corrigendum when using its statements.
\url{https://arxiv.org/abs/1401.8052}.
\end{thebibliography}
\end{document}
